\begin{lemma}[Scheide maximum-triple fact (h)]
Every color which is defective or free for a globally maximum triple
$(T,e_0,\varphi)$ is unused on the colored tree edges of $T$.
\end{lemma}
\begin{proof}
A free color is unused by the set-difference definition
\noderef{TreeFreeColors}, so the free-color set is disjoint from
\noderef{TreeUsedColors}.

Let $\alpha$ be defective.  The statement-level package
\noderef{ScheideGlobalMaximumOrder} says that $X=V(T)$ is closed.  Since an
$\alpha$-colored boundary edge identity exists, closure implies that
$\alpha$ is missing at no vertex of $X$.

If $\alpha$ were used on $T$, there would be a colored tree edge after the
distinguished first edge with color $\alpha$, by
\noderef{TreeUsedColors}.  The Tashkinov-tree clause in
\noderef{GlobalMaximumTriple} would then provide an earlier tree vertex at
which $\alpha$ is missing.  That contradicts the preceding paragraph.
Thus every defective color is also disjoint from the used-color set, which
proves both asserted disjointness statements.
\end{proof}
